%%%%%%%%%%%%%%%%%%%%%%%%%%%%%%%%%%%%%%%%%%%%%%%%%%%%%%%%%%%%%%%%%%%%%%%%%%%%%%
% CAPITOLO 3 --- APPROCCI AGILE AL PROJECT MANAGEMENT
%
% Principio ordinatore: le sezioni procedono per LIVELLO DI PRESCRITTIVITA',
% cioe' per quanto ciascun testo dice su che cosa fare concretamente.
%   3.2  valori          (Manifesto: due pagine, nessuna procedura)
%   3.4  metodi di team  (Scrum, Kanban: ruoli, eventi, artefatti)
%   3.5  scala           (SAFe, LeSS: strutture organizzative complete)
% Questo consente di dichiarare un criterio di selezione argomentabile (3.3) e
% permette al capitolo 4 di confrontare a parita' di livello.
%
% Convenzione di questo file: il testo visibile e' l'ossatura argomentativa
% (spesso riutilizzabile come frase di apertura della sezione); i blocchi
% commentati con % contengono le indicazioni di stesura e le fonti da usare.
% Decisioni prese: taglio generico di project management, senza pratiche
% ingegneristiche di dettaglio; Disciplined Agile rimandato al capitolo 4.
%%%%%%%%%%%%%%%%%%%%%%%%%%%%%%%%%%%%%%%%%%%%%%%%%%%%%%%%%%%%%%%%%%%%%%%%%%%%%%

\chapter{Approcci Agile al Project Management}

Il capitolo precedente ha potuto assumere due testi autoritativi come unica
fonte e derivarne l'intera trattazione. Per l'Agile un equivalente non esiste:
non c'\`e uno standard, non c'\`e un ente che lo definisca, e i testi
disponibili vanno da una dichiarazione di valori di due pagine a framework
proprietari che prescrivono l'intera struttura organizzativa. Il capitolo \`e
quindi ordinato secondo il \textbf{livello di prescrittivit\`a} dei testi
esaminati, dal meno al pi\`u vincolante.

% ===========================================================================

% ===========================================================================
\section{Il cambio di paradigma e le sue radici}

Gli approcci Agile rispondono a una classe di problemi per cui i metodi
``Tradizionali'' definiti precedentemente non sono stati costruiti, e lo fanno con
pratiche in uso da anni (prima del 2001).

\subsection{Il problema a cui l'Agile risponde}

Gli approcci descritti nel capitolo precedente presuppongono che l'esito atteso
sia specificabile con sufficiente precisione prima dell'inizio del lavoro. Nello
sviluppo software questa condizione \`e frequentemente disattesa.

L'\emph{Agile Practice Guide} del PMI nomina la distinzione sottostante: il
lavoro di progetto si colloca lungo un intervallo che va dal \emph{definable
work} al lavoro a elevata incertezza, o \emph{high-uncertainty work}. Il primo
\`e caratterizzato da procedure gi\`a rivelatesi efficaci su progetti simili,
con dominio e processi produttivi noti e livelli bassi di incertezza e di
rischio in esecuzione; il secondo comporta progettazione inedita e
soluzione di problemi, e richiede la collaborazione di specialisti
\cite{agilepracticeguide2}.

Il lavoro a elevata incertezza presenta tassi di cambiamento, complessit\`a e
rischio alti, e questo mette in difficolt\`a gli approcci predittivi, che
determinano in anticipo la maggior parte dei requisiti e governano le variazioni
attraverso un processo di richiesta di modifica.

Nel software concorrono due fattori. Il primo \`e l'assenza di vincoli fisici di
produzione: la modifica tardiva non impone di demolire quanto gi\`a realizzato,
e il suo costo dipende da come il prodotto \`e costruito pi\`u che da quanto
lavoro \`e stato svolto.
Il secondo \`e il momento in cui arriva la verifica: in un ciclo predittivo la conferma di
avere costruito la cosa giusta si ottiene alla consegna, cio\`e nel punto in cui
correggere costa di pi\`u.

La stessa guida del PMI registra l'esito di questa condizione: lo sviluppo
software commerciale degli anni novanta e dei primi anni duemila era gravato da
metodologie rigide, e i progetti venivano consegnati spesso fuori budget,
in ritardo e senza soddisfare le esigenze del committente.

Da qui la domanda a cui gli approcci Agile rispondono: come si governa un
progetto quando l'oggetto da costruire non \`e conoscibile in anticipo?

\subsection{Le radici: lo sviluppo iterativo prima del 2001}

Lo sviluppo iterativo e incrementale non nasce con il Manifesto Agile: lo
precede di decenni, ed \`e questa anteriorit\`a a spiegare perch\'e nel 2001
esistessero gi\`a diversi metodi da unificare.

Larman e Basili ne ricostruiscono l'uso a partire dagli anni cinquanta.
Il programma aeronautico X-15 gi\`a sfrutta principi assimilabili all'agile, e
parte del personale che lavorava al progetto passa poi al Project Mercury, che all'inizio
degli anni sessanta lavora per iterazioni di mezza giornata a durata fissata,
scrivendo le prove prima del codice. E fra il 1977 e il 1980 il software
dello Space Shuttle arriva alla delivery attraverso 17 iterazioni su 31 mesi.

I metodi che il Manifesto unificher\`a esistono tutti prima del 2001: Scrum
nasce all'inizio degli anni novanta alla Easel Corporation con iterazioni di
trenta giorni, DSDM viene definito nel gennaio 1994 da un gruppo di sedici
professionisti britannici, il \emph{Rational Unified Process} a met\`a decennio,
Extreme Programming dal 1996 sul progetto Chrysler C3 e Feature-Driven
Development nel 1997 \cite{larman2003}.

Il Manifesto non fonda dunque i metodi agili: li ratifica e fornisce loro un
lessico comune.

\section{Il Manifesto Agile}

\subsection{Genesi e statuto del testo}

Il \emph{Manifesto for Agile Software Development} \`e il risultato di un
incontro tenutosi dall'11 al 13 febbraio 2001 presso lo Snowbird ski resort, in
Utah, fra diciassette persone provenienti da metodi gi\`a esistenti ---
Extreme Programming, Scrum, DSDM, Adaptive Software Development, Crystal,
Feature-Driven Development, Pragmatic Programming e altri \cite{agilemanifesto}.

Il testo si compone di una dichiarazione di meno di settanta parole, seguita da
dodici principi.
Non \`e uno standard: nessun organismo ne rivendica la manutenzione e il testo
non ha subito revisioni dal 2001. Non \`e un metodo: non nomina attivit\`a,
ruoli o documenti, e non consente quindi di derivare che cosa fare. Non \`e
certificabile, non essendovi un ente che ne attesti la conformit\`a.

I due testi del capitolo precedente ricavano la propria autorit\`a da un
soggetto che li pubblica, li aggiorna e ne certifica l'applicazione. Il
Manifesto \`e firmato da diciassette persone a titolo individuale, e la sua
autorit\`a deriva unicamente dall'adesione successiva di chi lo ha assunto come
riferimento. L'assenza di un proprietario spiega tanto la rapidit\`a della
diffusione quanto l'appropriazione commerciale del termine, di cui si dir\`a
nella sezione~\ref{sec:critiche}.

\subsection{I quattro valori}

I quattro valori sono formulati come confronti fra coppie di termini, riportati
nella tabella~\ref{tab:valori-manifesto}.

\begin{table}[htbp]
\centering
\small
\caption{I quattro valori del Manifesto Agile, nella traduzione italiana
ufficiale.}
\label{tab:valori-manifesto}
\begin{tabularx}{\textwidth}{@{}XX@{}}
\toprule
\textbf{Si valorizza} & \textbf{pi\`u di} \\
\midrule
Gli individui e le interazioni & i processi e gli strumenti \\
Il software funzionante & la documentazione esaustiva \\
La collaborazione col cliente & la negoziazione dei contratti \\
Rispondere al cambiamento & seguire un piano \\
\bottomrule
\end{tabularx}
\end{table}

La clausola che chiude l'elenco \`e la parte del testo pi\`u frequentemente
omessa: ``fermo restando il valore delle voci a destra, consideriamo pi\`u
importanti le voci a sinistra'' \cite{agilemanifesto}.
Il Manifesto non prescrive di non documentare, di non adottare
processi, di non negoziare contratti o di non pianificare: stabilisce una
precedenza in caso di conflitto fra i termini.

\subsection{I dodici principi}

I dodici principi esplicitano i valori senza raggiungere il livello della
procedura. La tabella~\ref{tab:principi-manifesto} li raggruppa per famiglia,
indicando la numerazione originale.

\begin{table}[!htb]
\centering
\small
\caption{I dodici principi del Manifesto raggruppati per famiglia
\cite{agilemanifesto}.}
\label{tab:principi-manifesto}
\begin{tabularx}{\textwidth}{@{}>{\raggedright\arraybackslash}p{2.8cm}%
                             >{\raggedright\arraybackslash}p{1.8cm}X@{}}
\toprule
\textbf{Famiglia} & \textbf{Principi} & \textbf{Contenuto} \\
\midrule
Consegna e valore & 1, 3, 7 &
La soddisfazione del committente si persegue con rilasci frequenti e continui di
software funzionante, che costituisce anche la misura primaria
dell'avanzamento \\
\addlinespace
Cambiamento & 2 &
I requisiti mutevoli vanno accolti anche a sviluppo avanzato e trattati come
vantaggio competitivo per il committente \\
\addlinespace
Persone e organizzazione & 4, 5, 6, 8, 11 &
Collaborazione quotidiana fra committente e realizzatori; team costruiti attorno
a persone motivate, dotate dell'ambiente e della fiducia necessari;
conversazione diretta come mezzo di comunicazione pi\`u efficace; ritmo di
lavoro sostenibile a tempo indefinito; architetture e requisiti come esito di
team auto-organizzati \\
\addlinespace
Eccellenza tecnica & 9, 10 &
Attenzione continua alla qualit\`a tecnica e alla progettazione; semplicit\`a
intesa come arte di massimizzare il lavoro non svolto \\
\addlinespace
Riflessione e adattamento & 12 &
A intervalli regolari il team esamina il proprio funzionamento e corregge di
conseguenza il proprio comportamento \\
\bottomrule
\end{tabularx}
\end{table}
\newpage

Due principi meritano attenzione perch\'e vincolanti: il
sesto trova nella conversazione diretta il mezzo pi\`u efficace di trasmettere
informazione, e presuppone quindi team di dimensione ridotta e collocati nello
stesso luogo; l'undicesimo attribuisce a team auto-organizzati l'origine delle
architetture e dei requisiti, e sottrae dunque quelle decisioni a un livello
gerarchico superiore. Entrambi per\`o valgono solo per il singolo team: nessuno
dei dodici principi dice come gestire la pluralit\`a dei team, ed \`e su questa
omissione che nasce il problema della scala.

\subsection{Che cosa il Manifesto non dice}

Essendo una trattazione limitata, molti temi non vengono trattati e l'elenco delle
assenza risulta pi\`u istruttivo rispetto all'elenco dei suoi contenuti.

Il testo non nomina ruoli, riunioni ricorrenti, artefatti, stime, iterazioni
n\'e liste di lavoro arretrato. Tutto ci\`o che il senso comune associa
all'Agile proviene dai metodi, non dal Manifesto.

Il testo non nomina mai il project management. Parla di \emph{software
development}, e il termine \emph{project} vi compare solo come contenitore del
lavoro, nei principi quattro e cinque. Non si trova il progetto come
organizzazione temporanea, la giustificazione economica dell'investimento, la
governance, il rischio: cio\`e gli oggetti su cui sono costruiti per intero i
due testi precedenti. Il Manifesto e gli standard tradizionali non
sono quindi due risposte alternative alla stessa domanda, ma due testi che
parlano di cose diverse.

Il testo non dice infine nulla sulla scala, sui rapporti contrattuali e sugli
ambienti soggetti a vincoli normativi.

Proprio perch\'e non \`e operativo, il Manifesto richiede uno strato che lo
traduca in pratica. \`E la funzione che i metodi esaminati nel seguito
assolvono, ed \`e anche la ragione per cui sono molti.

\section{Dal valore alla pratica: il problema dell'operativizzazione}

Fra l'enunciato ``gliindividui e le interazioni pi\`u che i processi e gli strumenti''
e una decisione operativa esiste una distanza da colmare. \`E questa
distanza a spiegare la nascita di framework che caratterizza l'Agile e a
distinguerlo dagli approcci tradizionali.

L'\emph{Agile Practice Guide} descrive la distanza come una struttura a strati:
una mentalit\`a definita da quattro valori, guidata da dodici principi e resa
operativa da pratiche, fra le quali chi lavora sceglie in base alle proprie
necessit\`a \cite{agilepracticeguide2}.

Non \`e necessario adottare
alcuno degli approcci esistenti, perch\'e un approccio agile pu\`o essere
costruito da zero purch\'e aderisca alla mentalit\`a, ai valori e ai principi
del Manifesto.

A questo punto della trattazione è necessario definire alcuni termini:
\begin {itemize}
\item \emph{framework}: un sistema o una struttura di base di idee o
di fatti a sostegno di un approccio
\item  \emph{approccio agile} un approccio di sviluppo che applica mentalit\`a,
valori e principi del Manifesto attraverso pratiche
\end{itemize}
A questi si aggiunge la distinzione fra
metodo (che prescrive una sequenza di passi) e \emph{tool kit} (che raccoglie
opzioni fra cui scegliere senza prescriverne una).

Non essendo possibile raccogliere tutti gli approcci esistenti in un'unica trattazione,
la guida ne seleziona alcuni attraverso i seguenti criteri:
\begin {itemize}
\item concepiti per un uso olistico, cio\`e destinati a guidare un insieme ampio
di attivit\`a di progetto e non una singola attivit\`a
\item formalizzati per un uso comune, cio\`e non proprietari di una singola
organizzazione o di un singolo contesto
\item diffusi nell'uso corrente, secondo un insieme di indagini di settore
recenti
\end {itemize}

Gli approcci cos\`i selezionati sono raffigurati dalla mappa seguente dove
si distinguono i metodi di team dagli approcci scalati.

\begin{figure}[htbp]
\centering
\resizebox{\textwidth}{!}{%
\begin{tikzpicture}[font=\footnotesize,
  team/.style={draw, rounded corners=3pt, align=center, minimum height=0.6cm,
               inner xsep=4pt, fill=white},
  scal/.style={draw, rounded corners=3pt, align=center, minimum height=0.6cm,
               inner xsep=4pt, fill=black!12},
  sel/.style={line width=1.1pt}]

  % assi
  \draw[-{Latex[length=2.5mm]}] (-0.3,0) -- (12.2,0);
  \draw[-{Latex[length=2.5mm]}] (-0.3,0) -- (-0.3,7.6);
  \node[font=\small, anchor=north] at (6,-0.45)
    {Profondit\`a del dettaglio prescrittivo};
  \node[font=\small, rotate=90, anchor=south] at (-0.95,3.8)
    {Ampiezza della copertura del ciclo di vita};

  % approcci scalati
  \node[scal]      at (2.24,6.08) {Lean};
  \node[scal, sel] at (10.24,6.55) {SAFe};
  \node[scal, sel] at (3.34,5.03) {LeSS};
  \node[scal]      at (2.11,3.93) {Scrum of\\Scrums};

  % metodi di team
  \node[team]      at (5.76,5.54) {DSDM};
  \node[team]      at (7.83,4.11) {Agile UP};
  \node[team]      at (5.20,3.44) {Crystal\\Methods};
  \node[team]      at (5.20,2.40) {XP};
  \node[team, sel] at (1.73,2.21) {Kanban};
  \node[team]      at (1.79,1.30) {FDD};
  \node[team, sel] at (2.75,0.50) {Scrum};

  % legenda
  \node[team] (lt) at (10.4,2.3) {Metodo di team};
  \node[scal] (ls) at (10.4,1.5) {Approccio scalato};
  \node[draw, dashed, fit=(lt)(ls), inner sep=5pt] {};
\end{tikzpicture}}
\caption{Gli approcci agili disposti per ampiezza di copertura del ciclo di vita
e profondit\`a del dettaglio prescrittivo. Il bordo spesso segnala i quattro
approcci trattati per esteso nel seguito. Rielaborazione da
\cite{agilepracticeguide2}.}
\label{fig:mappa-approcci}
\end{figure}
\newpage

Assunta tale mappa come griglia, le due sezioni seguenti trattano per esteso
quattro casi, scelti per posizione e non per diffusione:
\begin {itemize}
\item Scrum, metodo di team a prescrittivit\`a media e ampiamente adottato
\item Kanban, metodo di team a prescrittivit\`a minima, che si applica sopra il
processo esistente invece di sostituirlo
\item SAFe e LeSS, due approcci scalati collocati agli estremi opposti dello
spettro prescrittivo
\end {itemize}

Gli altri sono richiamati in forma tabellare al termine di ciascuna sezione.

\newpage

\section{I metodi di team}

% Nota di taglio: la tesi resta su un piano di project management. Le pratiche
% tecniche vengono nominate per completezza dello strato, non descritte.

\subsection{Scrum}

% Da sviluppare:
% - Fondamento: il controllo empirico di processo, con i tre pilastri
%   trasparenza, ispezione, adattamento. E' il corrispettivo funzionale del
%   controllo per baseline degli approcci predittivi, ed e' su questo che il
%   capitolo 4 potra' impostare il confronto.
% - Le tre accountability (Product Owner, Scrum Master, Developers), i cinque
%   eventi (Sprint, Sprint Planning, Daily Scrum, Sprint Review, Sprint
%   Retrospective), i tre artefatti con i rispettivi commitment (Product
%   Backlog / Product Goal, Sprint Backlog / Sprint Goal, Increment /
%   Definition of Done).
% - Lo sprint come timebox e come unita' di pianificazione: si pianifica in
%   dettaglio solo l'iterazione corrente. Da collegare all'orizzonte di
%   pianificazione di PRINCE2 (2.3.6) e al rolling wave planning del PMBOK
%   (2.2.6): e' un'affinita' importante, ma il confronto va fatto nel cap. 4.
%
% MOSSA DIACRONICA, la stessa che ha funzionato sul PMBOK nel capitolo 2:
% la revisione 2020 rispetto alla 2017 rende il framework MENO prescrittivo,
% non piu'. Elimina il "development team" come entita' distinta per evitare
% dinamiche di delega interna al team; introduce il Product Goal; passa da
% self-organizing a self-managing (il team decide anche che cosa fare, non solo
% come); toglie le tre domande obbligatorie dal Daily Scrum; dichiara
% l'obiettivo di essere "minimally sufficient" e scende sotto le tredici
% pagine. Un framework che nel tempo rimuove prescrizioni e' un dato in se',
% e va confrontato con la traiettoria opposta del PMBOK, che nell'ottava
% edizione reintroduce i processi (di nuovo: osservazione da cap. 4).
% - Segnalare anche cio' che Scrum NON copre: nessuna pratica tecnica, nessuna
%   indicazione su architettura, contratti, budget, portfolio. E' il vuoto che
%   XP occupa (3.4.3) e che la scala deve colmare (3.5).

\subsection{Kanban}

% Da sviluppare, fonte: l'annex A3.4 dell'Agile Practice Guide
% \cite{agilepracticeguide2} per il metodo, e il par. 2.3 della stessa guida
% per i concetti lean da cui discende.
% - Le pratiche fondamentali: visualizzare il flusso di lavoro, limitare il
%   lavoro in corso (WIP), gestire il flusso, rendere esplicite le politiche,
%   istituire cicli di feedback, migliorare in modo collaborativo ed
%   evolutivo.
% - Il principio di partenza: "start with what you do now". Kanban e' un metodo
%   di cambiamento evolutivo, non un modello organizzativo da adottare. Questa
%   e' la differenza che conta per la tesi: non richiede di riorganizzare il
%   lavoro prima di cominciare, il che ne spiega l'adozione in contesti dove
%   Scrum non e' applicabile (manutenzione, supporto, flussi non progettuali).
% - La schedulazione a flusso e la teoria dei vincoli: il PMBOK 8 le tratta
%   esplicitamente \cite{pmbok8}, quindi qui esiste gia' un punto di contatto
%   documentato fra i due mondi. Segnalarlo senza svilupparlo.
% - Metriche proprie: lead time, cycle time, throughput, invece di velocity.
% Conclusione: Kanban dimostra che "agile" non implica "iterativo a timebox",
% e quindi che l'opposizione predittivo/iterativo non esaurisce la questione.

\subsection{Le pratiche ingegneristiche e il loro statuto}

% Trattazione volutamente breve: la tesi resta sul piano del project
% management e non entra nel merito delle singole pratiche.
% Da dire, e basta:
% - XP prescrive pratiche di realizzazione oltre che di organizzazione del
%   lavoro, ed e' l'unico dei metodi qui considerati a farlo.
% - Il punto rilevante per il project management non e' quali siano quelle
%   pratiche, ma la TESI che XP sostiene: che l'adattabilita' del processo
%   dipenda da proprieta' tecniche del prodotto, e che nessuna riorganizzazione
%   del lavoro possa produrre agilita' se il costo della modifica resta alto.
%   E' un'affermazione forte, perche' colloca fuori dalla portata del project
%   manager una delle condizioni del successo del progetto.
% - Da qui una domanda che il capitolo 4 puo' riprendere: se cio' e' vero,
%   il perimetro tradizionale del project management e' sufficiente?
% Non introdurre esempi di pratiche specifiche.

\subsection{Gli altri metodi di team}

% Trattazione tabellare, non discorsiva. Una tabella con: metodo, anno,
% autore/i, elemento distintivo, e una riga ciascuno per Crystal, FDD, DSDM e
% Adaptive Software Development. Servono a documentare che il campo e' piu'
% ampio dei due casi trattati e a giustificare, per contrasto, la selezione
% dichiarata in 3.3. Fonte per la collocazione relativa: la mappa del
% PMBOK 8 \cite{pmbok8}.

\section{La scala}

\subsection{Perch\'e la scala \`e un problema}

% Da sviluppare, ed e' la sottosezione concettualmente piu' densa di 3.5:
% - Il Manifesto presuppone team piccoli, co-locati, in comunicazione diretta.
%   Il principio sulla conversazione faccia a faccia e quello sui team
%   auto-organizzati non si estendono per semplice moltiplicazione.
% - Coordinare piu' team richiede sincronizzazione, allineamento
%   architetturale, gestione delle dipendenze e allocazione di budget: cioe'
%   esattamente le funzioni che il capitolo 2 ha visto attribuite alla
%   governance. La scala reintroduce quindi il problema che il livello di team
%   aveva potuto ignorare.
% - Da qui l'obiezione ricorrente, da riportare in modo equanime: i framework
%   di scala reintrodurrebbero cio' che l'Agile aveva rimosso, e sarebbero un
%   ritorno del predittivo sotto lessico agile. Presentare l'obiezione e la
%   replica, senza prendere posizione qui (la posizione va nel cap. 4).

\subsection{SAFe}

% Da trattare come l'estremo prescrittivo dello spettro.
% Fonte: l'annex A3.8.2 dell'Agile Practice Guide \cite{agilepracticeguide2}.
% Il framework e' proprietario e a revisione frequente: descriverlo attraverso
% la guida evita di doverne inseguire le versioni.
% - Struttura a livelli (team, programma/ART, portfolio), il Program Increment
%   e il PI planning come evento di sincronizzazione, i ruoli aggiuntivi.
% - Il dato interessante per la tesi: SAFe e' il framework agile piu' adottato
%   nelle grandi organizzazioni ed e' anche il piu' criticato dall'interno del
%   movimento. Le due cose insieme dicono qualcosa sul rapporto fra
%   adottabilita' e fedelta' ai principi, che il capitolo 5 potra' verificare
%   empiricamente.
% - Non inseguire le versioni: descrivere l'impianto, non il dettaglio
%   corrente.

\subsection{LeSS}

% L'estremo opposto: fonte l'annex A3.8.3 dell'Agile Practice Guide
% \cite{agilepracticeguide2}.
% - L'idea di "descaling": non aggiungere strutture di coordinamento ma
%   rimuovere complessita' organizzativa; un solo Product Backlog, un solo
%   Product Owner, uno Sprint comune, feature team invece di component team.
% - Il contrasto con SAFe e' il vero contenuto della sottosezione: due risposte
%   opposte allo stesso problema. Vale la pena impostarlo come coppia, perche'
%   mostra che "scalare l'Agile" non e' una questione tecnica ma una scelta
%   su quanta struttura sia accettabile.

\subsection{Gli altri approcci scalati}

% Solo citati, breve tabellare

\section{Che cosa l'Agile dice del project management}

% SEZIONE CARDINE DEL CAPITOLO. E' quella che lo rende un capitolo di project
% management e non di ingegneria del software, ed e' l'aggancio diretto al
% filo conduttore della tesi (l'interpretazione del Project Management dal
% punto di vista dei diversi approcci).

\subsection{L'assenza del project manager}

% Da sviluppare:
% - Constatazione documentale: verificarla e dichiararla per ciascuna fonte
%   (Manifesto, Scrum Guide, Kanban), perche' e' un'osservazione forte e va
%   sostenuta.
% - L'assenza non e' una svista ma una conseguenza: se il team e'
%   auto-gestito, il coordinamento interno non richiede un ruolo dedicato.
% - Da segnalare che il PMBOK 8 registra questa situazione dal proprio lato,
%   osservando che negli approcci adattivi funzioni di project management
%   possono essere assorbite da product owner, scrum master o agile coach, e
%   che "l'essenza del project management risiede nelle caratteristiche del
%   progetto stesso piuttosto che nel titolo dell'individuo che lo
%   supervisiona" \cite{pmbok8}. Qui il richiamo e' legittimo perche' la
%   sezione parla del rapporto fra Agile e project management in generale.

\subsection{La ridistribuzione delle funzioni di project management}

La domanda corretta non \`e se l'Agile abbia un project manager, ma dove
finiscano le funzioni che il capitolo precedente ha visto attribuite a quel
ruolo.

% Da sviluppare, ed e' il contributo analitico piu' originale del capitolo.
% Metodo suggerito: prendere le funzioni identificate nel capitolo 2 e
% chiedersi per ciascuna dove si collochi in un impianto agile. Una tabella a
% due colonne (funzione / collocazione) sarebbe efficace. Funzioni da
% considerare:
%   - giustificazione dell'investimento e sua rivalidazione;
%   - definizione e controllo dell'ambito;
%   - pianificazione e stima;
%   - controllo dell'avanzamento e reportistica;
%   - gestione degli stakeholder;
%   - gestione del rischio;
%   - gestione della qualita';
%   - chiusura e realizzazione dei benefici.
% Esiti attesi, da verificare fonte alla mano: alcune funzioni si spostano sul
% Product Owner (valore, priorita', ambito), altre sul team (stima,
% pianificazione di dettaglio, qualita'), altre risalgono all'organizzazione
% (finanziamento, benefici), altre ancora NON TROVANO COLLOCAZIONE ESPLICITA
% in nessuna delle fonti agili esaminate: il rischio e la giustificazione
% continuativa dell'investimento sono i candidati principali.
% Quest'ultimo punto e' importante e va detto senza indulgenza: e' una lacuna
% documentale dei testi agili, non un difetto delle organizzazioni che li
% adottano, ed e' una delle ragioni per cui nascono gli approcci ibridi.

\subsection{Progetto o prodotto}

% Sottosezione breve ma concettualmente decisiva.
% - L'Agile ragiona per PRODOTTO: flusso continuo di valore, team stabili nel
%   tempo, backlog che non si esaurisce. Il progetto e' per definizione
%   TEMPORANEO, e la temporaneita' e' nella definizione di progetto di tutte le
%   fonti citate nel capitolo 2, ISO compresa \cite{iso21500-2021}.
% - Ne segue una tensione concettuale reale e non risolta: un'organizzazione
%   che adotta pienamente il modello a prodotto smette in un certo senso di
%   fare progetti, e con essi perde i punti di decisione, la chiusura e la
%   rendicontazione finale.
% - Questa tensione va lasciata aperta qui: e' materiale per la sezione
%   "Verso una definizione sintetica di Project Management" del capitolo 4.

\section{Critiche, limiti ed evidenza empirica}
\label{sec:critiche}

